% !TEX root = ../main.tex
\chapter{앞선 연구 살펴보기와 이론의 틀}
\label{ch:literature}

\ref{ch:introduction}장에서 세운 연구 문제는 다음과 같아요. 송금(transfer) 그래프(점과 화살표로 이은 그림)로 매기는 평판 점수는, 토큰을 대신 써도 된다고 맡기는 일에 아무 역할도 주지 않아요. 이 맡김은 믿음을 분명하게 드러내는 행동이고, 모든 ERC-20(토큰 공통 규칙) 토큰이 이것을 기록해요. 또 이런 점수의 그래프는 돈을 낼 때마다 커져요. 목표 O1--O5와 연구 질문 RQ1--RQ5는 이 문제에서 나와요. 이 연구에서는 허락(allowance, 토큰을 얼마까지 대신 써도 된다는 허락)과 송금을 함께 따라가는 페이지랭크(PageRank) 걷기(화살표를 따라 점에서 점으로 옮겨 다니는 것)를 살펴봐요. 그리고 이 걷기를 허락과 송금 가운데 하나씩만 따라가는 두 걷기와 견주어요. 바로 엔도스랭크(EndorseRank)와 AWP(적응형 가중 페이지랭크)예요. 이를 위해 각 방법의 그래프에 페이지랭크를 돌리고, 그렇게 나온 순위를 순위 상관(두 줄 세우기가 얼마나 비슷한지 나타내는 값)으로 점검해요. 아직 남은 빈틈(밝혀지지 않은 부분)은 \ref{ch:methodology}장과 \ref{ch:results}장에서 다뤄요.

\dissertationSubheading[sec:blockchain-defi-foundations]{블록체인과 디파이의 기초}

누구나 볼 수 있는 공개 블록체인은 거래를 해시(내용을 짧은 지문 같은 값으로 바꾼 것)로 묶은 블록에 기록해요. 그리고 각 블록은 그 뒤에 덧붙는 블록들이 지켜 줘요 \parencite{nakamoto2008bitcoin}. 참여하는 컴퓨터(블록체인 노드) 가운데 절반이 넘는 수가 한 블록을 받아들이고 나면, 그 블록을 고쳐 쓰는 데 큰 비용이 들어요. 이 비용 덕분에, 거래하는 사람들은 저마다 알려진 법적 신분이 없어도 장부 위에서 거래를 마무리(결제)할 수 있어요(\emph{믿음을 최소로 줄인} 결제).

비트코인(Bitcoin)은 이런 장부 위에 사람과 사람이 직접 주고받는(P2P) 전자 현금 네트워크를 만들었어요 \parencite{nakamoto2008bitcoin}. 이더리움(Ethereum)도 뒤에 덧붙이기만 할 수 있는 장부를 그대로 써요. 하지만 계정이 잔액과 함께 프로그램도 가질 수 있게 했어요. 그러면 거래 하나가 그 프로그램의 함수(정해진 일을 하는 명령 묶음) 하나를 불러서 계약의 상태를 바꿀 수 있어요 \parencite{buterin2014ethereum}. 이 함수들은 이더리움 가상 머신(EVM, 이더리움 안에서 프로그램을 돌리는 장치)에서 돌아가요 \parencite{wood2025yellowpaper}. 이렇게 장부에 저장된 프로그램을 스마트 계약이라고 불러요 \parencite{szabo1997formalizing}. 정직한 노드라면 어느 것이든, 같은 함수 호출을 다시 돌려 보면 똑같은 결과를 얻어요 \parencite{wood2025yellowpaper}. 그 결과, 담보(돈을 빌릴 때 맡겨 두는 재산) 비율, 청산(손실이 너무 커져 거래가 강제로 끝나는 일) 조건, 수수료 표를 믿을 만한 제삼자 대신 계약이 지키게 할 수 있어요 \parencite{schar2021defi}.

이런 계약들이 모두 모여 디파이(DeFi, 탈중앙 금융)를 이루어요. 이 계약들은 함께 거래소, 대출 플랫폼, 파생상품(다른 자산의 값에 따라 값이 정해지는 상품) 시장, 자산 관리자가 하는 일을 해요. \textcite{werner2022sok}에서는 서비스(프로토콜)의 종류를 정의했어요. 그리고 \textcite{schar2021defi}에서는 서비스들이 서로의 위에 어떻게 층층이 쌓이는지 설명했어요. 한 서비스가 다른 서비스와 맞물려 움직일 수 있어서, 이미 있는 금융 상품들을 엮어 새 금융 상품을 만들 수 있어요 \parencite{kitzler2023defi}. 그런데 바로 이 성질 때문에 한 곳의 고장이 여러 서비스로 번질 수도 있어요. 실제로 성공한 공격에는 플래시 론(한 거래 안에서 빌리고 바로 갚는 대출)을 악용한 공격이 있었어요 \parencite{qin2021flashloans}. 가격 오라클(블록체인 밖의 가격을 알려 주는 장치)을 조작한 공격도 있었어요 \parencite{zhou2023sok}. 이런 블록체인의 거래 기록은 누구나 볼 수 있어요. 하지만 그 안의 주소들은 가명의(진짜 이름 대신 주소로만 활동하는) 주소예요. 그래서 보통 금융에서 하는 신분 확인을 하기가 어려워요 \parencite{lin2025riskprop}.

\dissertationSubsubheading[sub:pseudonymity-credit]{가명성과 신용 문제}

보통의 신용 점수 매기기는 신용 기록을 법적 주체(법으로 인정받는 사람이나 회사)에 묶어요. 그런데 공개 블록체인에서는 분석의 기본 단위가 주소(지갑)예요. 주소 하나가 나타내는 것은 법적 주체일 수도, 신탁(남의 재산을 맡아 관리하는 단체)일 수도, 알고리즘(정해진 계산 순서)으로 움직이는 대리인일 수도, 스마트 계약일 수도 있어요. 한 주체가 주소를 여러 개 가질 수 있어요 \parencite{douceur2002sybil}. 그리고 통계로 주소를 묶어 보는 방법(클러스터링)으로도 누가 주인인지는 일부만 알아낼 수 있어요 \parencite{victor2020clustering}. 그래서 온체인(블록체인 위의) 대출에서는 빌리는 사람이 빌린 돈보다 값이 더 나가는 담보를 맡겨야 해요. 담보의 값이 정해 둔 선 아래로 떨어지면, 청산하는 사람이 빚의 일부를 대신 갚고 그만큼의 담보를 싸게 받아 갈 수 있어요 \parencite{bartoletti2021lending}. 이 방식은 안전하지만, 돈을 묶어 둬요. 온체인 신용 점수는 블록체인 위의 행동만 보고, 어떤 참여자가 돈을 갚을 가능성이 얼마나 되는지를 수로 나타내려고 해요. 최근의 제안들은 전망 이론(사람이 이익과 손해를 어떻게 느끼는지 설명하는 이론)을 바탕으로 하기도 해요 \parencite{hassija2020lending}. 대출 자료에 쓴 기계 학습을 바탕으로 하기도 해요 \parencite{kotb2023credit,palaiokrassas2024credit}. 그리고 계정 연결망을 따라 위험이 퍼지는 모형을 바탕으로 하기도 해요 \parencite{lin2025riskprop}.

\dissertationSubheading[sec:ethereum-l2]{이더리움, 레이어 2 롤업, 아비트럼 원}

대부분의 디파이 앱은 이더리움 메인넷(이더리움의 본래 블록체인)에 올라갔어요. 이더리움 메인넷은 예전부터 디파이의 새로운 시도가 일어나고 유동성(시장에서 바로 사고팔 수 있는 돈)이 몰리는 주된 곳이었어요. 하지만 이더리움이 작업 증명(계산 문제를 풀어 블록을 만드는 방식)으로 돌아가던 동안에는, 안전을 약하게 만들지 않고서는 블록 크기를 크게 늘리거나 블록 사이 간격을 크게 줄일 수 없었어요 \parencite{gervais2016pow}. 그래서 메인넷을 쓰는 사람이 늘수록 가스(거래 수수료) 비용이 오르고, 처리량(정해진 시간에 처리할 수 있는 거래의 양)도 늘리기 어려웠어요. 레이어 2(L2) 서비스들은 이더리움의 안전 보장은 지키면서 거래를 더 싸게 만들려고 해요. 어떤 것은 거래를 오프체인(블록체인 밖의) 공간에서 실행해요. 다른 것은 거래를 묶어서, 실행 결과를 짧게 줄인 요약을 이더리움 메인넷에 올려요 \parencite{gudgeon2020layertwo}. \textcite{thibault2022rollups}에서는 롤업(rollup, 거래를 묶어 요약을 메인넷에 올리는 L2)을 두루 살펴보고, 그 구현(실제로 만든 방식)을 견주었어요. 아비트럼(Arbitrum) 같은 낙관적 L2(일단 맞다고 보고, 문제가 있을 때만 따지는 L2)는 제출된 거래 묶음을, 이의 제기 기간 안에 반대 증거가 나오지 않으면 올바른 것으로 봐요 \parencite{kalodner2018arbitrum}.

아비트럼 원(Arbitrum One)은 이런 낙관적 L2 가운데 하나이고, EVM에 바탕을 둔 스마트 계약을 실행할 수 있어요. 아비트럼 원은 디파이 앱들이 즐겨 쓰는 플랫폼이 되었어요. 그 가운데에는 무기한 선물 거래(끝나는 날 없이, 값이 오를지 내릴지에 거는 거래)를 하는 거래소인 GMX도 있어요 \parencite{gmxnddocs}. 이 연구에서 아비트럼 원을 쓰는 까닭은 두 가지예요. 공개된 빅쿼리(BigQuery) 자료에 아비트럼 원의 ERC-20 \texttt{Approval} 이벤트와 \texttt{Transfer} 이벤트(계약이 남기는 기록)가 모두 들어 있어요. 또 수수료가 충분히 낮아서, 살펴본 기간들 안에 활동이 빽빽하게 들어 있어요. 더 검증하지 않고서는, 우리 결과가 다른 L2나 이더리움 메인넷에도 똑같이 들어맞는다고 주장하지 않아요. 이 틀은 같은 로그를 남기는 블록체인이라면 어디에나 쓸 수 있어요.

\dissertationSubheading[sec:erc20-allowance-lifecycle]{ERC-20 토큰과 허락의 한살이}

ERC-20 표준은 대체 가능한 토큰(하나하나가 똑같은 값어치여서 서로 바꿔도 되는 토큰)을 위한 공통 인터페이스(프로그램끼리 서로 약속한 함수 모음)를 정해요. 이 인터페이스는 \texttt{transfer}, \texttt{approve}, \texttt{transferFrom}, \texttt{allowance}와 여기에 딸린 이벤트들로 이루어져요 \parencite{vogelsteller2015erc20}. 주인(허락을 해 주는 쪽)이 \texttt{approve(spender, amount)}를 부르면, 토큰 계약은 스펜더(spender, 허락을 받아 대신 쓰는 쪽)가 주인의 잔액에서 \texttt{amount}까지 옮겨도 된다고 기록해요. 그러면 스펜더는 \texttt{transferFrom}을 불러 토큰 송금을 실제로 할 수 있어요. EIP-2612의 \texttt{permit}(서명 허락)을 쓰면, 오프체인 서명 하나로 허락을 바꿀 수 있어요. 이 서명은 거래 한 번으로 처리돼요 \parencite{lundfall2020permit}. 허락이 뜻하는 바는 바뀌지 않아요.

그림~\ref{fig:allowance-sequence}에서는 ERC-20 허락의 한살이(생겨서 쓰이고 거두어지기까지)를 따라가 봐요. 허락은 쓸 권한을 주는 것이고, 실제로 토큰을 옮기지는 않아요. 토큰 주인은 라우터(거래를 이어 주는 계약), 금고(vault, 토큰을 맡아 두는 계약), 거래 계약에 허락을 줄 수 있어요. 정해 둔 한도 안에서 일할 것이라고 믿는 계약들이에요. 허락 거두기(0을 허락하기)와 바꾸기(0이 아닌 새 허락을 정하기)는 둘 다 지금 효력이 있는 허락을 바꿔요. 그리고 \texttt{transferFrom}은 그 허락을 줄여요. 어떤 토큰은 남은 양을 새 \texttt{Approval} 이벤트로 기록해요. 엔도스랭크는 주인, 스펜더, 토큰마다 가장 최근 허락 하나만 남겨요. 그래서 새 허락은 앞선 허락을 통째로 대신해요. 그리고 살아남은 허락은 효력이 그대로 있는 동안에도, 오래될수록 무게(가중치)가 줄어요(\ref{sec:allowance-endorsement}절).

\begin{figure}[htbp]
\centering
\begin{tikzpicture}[
  font=\small,
  actor/.style={draw, rounded corners, minimum width=2.4cm, minimum height=0.75cm, align=center},
  msg/.style={-{Stealth[length=2.5mm]}, thick},
  note/.style={font=\scriptsize\itshape, text=black!70, align=left}
]
% Actors
\node[actor] (owner) at (0,0) {주인};
\node[actor] (token) at (5.5,0) {ERC-20 토큰};
\node[actor] (spender) at (11,0) {스펜더};

% Lifelines (vertical dashed lines)
\draw[dashed, gray] (owner.south) -- ++(0,-6.6);
\draw[dashed, gray] (token.south) -- ++(0,-6.6);
\draw[dashed, gray] (spender.south) -- ++(0,-6.6);

% Message 1: approve (owner -> token), horizontal span 5.5cm
\draw[msg] (0,-1.2) -- node[above,font=\scriptsize] {\texttt{approve(spender,$a$)}} (5.5,-1.2);
\node[note] at (2.75,-1.7) {\texttt{Approval} 남기기; 허락 저장};

% Alternative to message 1: EIP-2612 permit, submitted by the spender (or a relayer) with the owner's signature
\draw[msg, densely dotted] (11,-2.6) -- node[above,font=\scriptsize] {\texttt{permit(owner,spender,$a$,sig)}} (5.5,-2.6);
\node[note] at (8.25,-3.1) {\texttt{approve} 대신: 주인이 오프체인에서 서명;\\ \texttt{Approval} 남기기; 허락 저장};

% Message 2: transferFrom (spender -> token), horizontal span 5.5cm
\draw[msg] (11,-4.2) -- node[above,font=\scriptsize] {\texttt{transferFrom(owner,to,$v\leq a$)}} (5.5,-4.2);
\node[note] at (8.25,-4.7) {토큰 옮기기; 허락 줄이기};

% Message 3: revoke (owner -> token), dashed
\draw[msg, dashed] (0,-5.8) -- node[above,font=\scriptsize] {\texttt{approve(spender,$0$)} (허락 거두기)} (5.5,-5.8);
\end{tikzpicture}
\caption[{ERC-20 허락의 한살이예요. 허락하기(approve) 또는 서명 허락(permit), 해도 되고 안 해도 되는 \texttt{transferFrom}, 허락 거두기로 이어져요.}]{ERC-20 허락의 한살이예요. \texttt{approve}나, 서명으로 대신하는 \texttt{permit}(점선), 해도 되고 안 해도 되는 \texttt{transferFrom}, 그리고 허락 거두기로 이어져요. 엔도스랭크는 점수를 매기는 기간 안에서 주인, 스펜더, 토큰마다 가장 최근 \texttt{Approval} 이벤트를 남겨요. 그리고 그 허락에 나이(지난 시간)와 달러(USD) 금액으로 무게를 매겨요. 나중 허락이 앞선 허락을 대신해요.}
\label{fig:allowance-sequence}
\end{figure}

가지고 있는 자료에서 허락은 송금에 비해 듬성듬성해요. 허락 한 번이 여러 번의 송금을 감당할 수 있어요. 그리고 많은 지갑은 바깥의 어떤 계정에도 허락을 한 번도 내준 적이 없어요. 이 희소성(화살표가 적음)은 반복(한 바퀴) 한 번에 드는 비용을 낮춰요. 또 이 희소성은 뜻의 차이도 보여 줘요. 엔도스랭크는 써도 된다는 권한을 수로 나타내지만, 송금 그래프는 돈이 흘러간 기록을 담기 때문이에요.

\dissertationSubheading[sec:graph-reputation]{그래프로 매기는 온체인 평판}

페이지랭크와 그 변형들은, 방향이 있는 화살표(엣지)를 따라 점(노드)으로 흘러 들어오는 영향력을 보고 점에 순위를 매겨요 \parencite{page1999pagerank}. 최근 두 연구가 이 방법을 블록체인 거래 그래프에 썼어요. \textcite{do2023hodgerank}에서는 사회적 신용(다른 사람들과의 관계에서 생기는 믿음)을 재려고 이더리움 거래에 페이지랭크와 HodgeRank를 돌렸어요. 이 점수는 금융 기관에 신용 기록이 없는 사람들을 위한 것이에요. \textcite{do2023awp}에서는 적응형 가중 페이지랭크(AWP)를 제안했어요. AWP는 이 논문이 엔도스랭크와 나란히 평가하는 송금 그래프 방법이에요(\ref{sec:awp-in-depth}절). \textcite{lin2025riskprop}에서는 이더리움 계정의 위험 등급 매기기를 점에 순위를 매기는 문제로 다뤘어요. 그리고 거래 연결망을 따라 위험을 퍼뜨려서(전파해서), 라벨(미리 알고 있는 정답 표시)이 없는 계정에도 등급을 매겼어요. \textcite{wu2022phishers}에서는 거래 그래프의 연결망 임베딩(그래프 속 위치를 숫자 목록으로 바꾼 것)으로 피싱(속여서 돈이나 정보를 빼앗는 사기) 주소를 찾아냈어요. 이 임베딩은 거래의 금액과 시간을 함께 따져요. \textcite{bellini2020trust}에서는 블록체인에 바탕을 둔 믿음과 평판 관리를 두루 살펴보았어요. \textcite{almasoud2020reputation}에서는 스마트 계약 위에 만든 평판 시스템들을 검토했어요. 이런 시스템에서는 사용자들이 온라인 서비스나 상품을 내놓는 쪽이 얼마나 믿을 만한지 점수를 매겨요. 암호화폐 거래를 대상으로 한 데이터 마이닝(많은 자료에서 쓸모 있는 규칙을 찾는 일)도 마찬가지로 거래 기록에서 출발해요. 주소를 서로 잇는 일과 사용자 행동 연구에 쓰여요 \parencite{liu2021knowledge}. \textcite{packin2024credit}에서는, 탈중앙 신용 점수가 치우친 결정을 바깥 사람이 들여다볼 수 없는 자동 처리 과정으로 바꿔 버릴 수 있다고 경고해요.

연구 문헌 밖에서, 실제로 쓰이는 디파이와 web3(블록체인 위에서 돌아가는 인터넷 서비스) 시스템은 보통 다른 방법으로 믿음을 쌓아요. 주소는 누구나 장부에서 살펴볼 수 있는 활동으로 평판을 쌓고, 그 주소에 묶인 도구가 그 활동을 요약해요. 그런 도구 가운데 하나가 휴먼 패스포트(Human Passport)예요. 예전 이름은 깃코인 패스포트(Gitcoin Passport)였어요. 이 도구는 스탬프(stamp)라고 부르는, 확인할 수 있는 자격 증명을 모아요. 스탬프는 정부 신분증, 생체 정보(얼굴이나 지문 같은 몸의 정보), 다른 사람이 대신 확인해 주기, 예전의 web2와 web3 활동 같은 점검마다 있어요. 스탬프마다 무게가 있는데, 설명서는 이 무게를 그 스탬프가 주는 사람이라는 신호의 세기라고 설명해요. 그리고 무게를 매겨 더한 값이 기준값에 이르면 그 주소는 통과해요 \parencite{humanpassportnddocs}. 같은 설명서에 따르면, 패스포트는 아비트럼을 포함한 여러 네트워크의 EVM 주소를, 온체인 활동을 보고 사람인지 시빌(한 사람이 만든 가짜 주소)인지로 나누기도 해요 \parencite{humanpassportnddocs}. 패스포트의 목적은 시빌 공격을 막는 것이에요. 그래서 온체인 증거와 오프체인 증거를 합치고, 등급을 매겨 줄 세우는 것이 아니라 들여보낼지 말지 정하는 문처럼 일해요. 사람 증명(주소 뒤에 진짜 한 사람이 있음을 보이는 것) 서비스들을 살펴본 한 검토 연구는, 이 서비스들이 시빌 공격을 어떻게 막는지 설명해요. 중앙 발급 기관의 증명 대신, 다른 사람이 대신 확인해 주기나 투표 같은 주관적인 입력을 쓴다는 것이에요 \parencite{siddarth2020watchmen}. 그리고 블록체인에 대한 시빌 공격을 살펴본 최근 검토 연구는, 아비트럼 같은 레이어 2 네트워크에서 일어난 에어드롭 파밍(공짜 토큰을 노린 가짜 활동)을 따라가 보고, 방어 수단 가운데 하나로 사람 증명 시스템을 평가해요 \parencite{bordeianu2026sybil}. 이런 도구는 한 주소를 따로 있는 한 사람이 다루는지에 답해요. 엔도스랭크는 다른 질문에 답해요. 다른 주소들이 자기 토큰을 쓰도록 허락해 주는 주소가 어느 것인지예요. 그리고 엔도스랭크는 온체인 기록만으로 그 주소들에 순서를 매겨요(표~\ref{tab:method-comparison-lit}).

그래프 방법이 매력 있는 까닭은 이래요. (i) 공개된 정보만 써요. (ii) 순위 매기기에 알맞은 연속적인 점수(끊김 없이 이어지는 값)를 내요. (iii) 화살표 수로 계산 복잡도(계산하는 데 드는 수고가 얼마나 빨리 늘어나는지)를 분석할 수 있어요. 하지만 그래프 방법은 시빌 공격에 약해요 \parencite{douceur2002sybil}. 자전 거래(자기끼리 사고팔아 거래가 많은 척하기)에도 약해요. 또 그래프 중심성(그래프에서 얼마나 가운데에 있는지)이 실제로 빚을 못 갚을 위험과 어긋나는 문제도 있어요. 페이지랭크 점수는 금융 규제에서 말하는 신용 점수가 아니에요(\ref{sub:non-claims}절).

퍼뜨린 점수는 더 단순한 점수도 이겨야 해요. 웹 그래프처럼 차수 상관(이어진 점들끼리 화살표 수가 서로 닮은 정도)이 약한 곳에서는, 페이지랭크가 평균적으로 들어오는 화살표 수(입차수)에 비례해요. 그래서 입차수로 페이지랭크를 어림할 수 있어요 \parencite{fortunato2008indegree}. 새 링크는 이미 링크를 많이 가진 점에 몰려요. 우선 연결(이미 많이 이어진 점에 새 링크가 더 잘 붙는다는 생각)이 이렇게 가정해요 \parencite{barabasi1999scaling}. 실제로 잰 연결망의 자람도 이것을 보여 줘요 \parencite{newman2001clustering}. 그래서 링크 예측(앞으로 어떤 링크가 생길지 미리 맞히기)에서는, 이더리움 거래 연결망에서도 \parencite{lin2022tracking}, 차수에 바탕을 둔 예측 방법을 표준 기준선(비교 기준)으로 써요 \parencite{libennowell2007link}.

\dissertationSubheading[sec:pagerank-formalism]{페이지랭크의 수식}

점이 $|V|$개인 방향이 있는 그래프 $G=(V,E)$를 생각해 봐요. $i$에서 $j$로 가는 화살표의 무게는 $w_{ij}>0$으로 써요. 이 무게는 $i$가 $j$에게 주는 영향력을 나타내요. 그 영향력은 앞에서 소개한 보증의 뜻이어도 되고, 송금의 뜻이어도 돼요. 이에 맞는 무게 인접 행렬(어느 점이 어느 점과 얼마의 무게로 이어졌는지 적은 표)을 $W$라고 해요. 그리고 점 $i$의 나가는 무게의 합을 $o_i=\sum_j w_{ij}$라고 해요. 그러면 줄마다 합이 1이 되는 이동 행렬(전이 행렬) $P$는 다음과 같아요.
\begin{equation}
P_{ij} =
\begin{cases}
w_{ij}/o_i & \text{만약 } o_i>0 \text{이면},\\
1/|V| & \text{만약 } o_i=0 \text{이면}
\end{cases}
\label{eq:transition}
\end{equation}
(매달린 점(나가는 화살표가 없는 점)은 자기 몫을 모든 점에 고르게 나눠 줘요). 감쇠 계수(화살표를 따라갈 확률) $d\in(0,1)$ 값이 주어지면, 구글 행렬(Google matrix)은 다음과 같아요.
\begin{equation}
G_{\mathrm{PR}} = d\, P^{\top} + \frac{1-d}{|V|}\,\mathbf{1}\mathbf{1}^{\top},
\label{eq:google-matrix}
\end{equation}
그리고 페이지랭크 벡터(수를 한 줄로 늘어놓은 것) $\mathbf{r}$은 $\mathbf{r}\geq 0$과 $\sum_i r_i=1$을 지키면서 $\mathbf{r}=G_{\mathrm{PR}}\mathbf{r}$을 만족해요 \parencite{page1999pagerank,langville2006pagerank}. $P$에서 매달린 점의 줄을 가운데 항으로 따로 떼어 쓰고 원소(벡터 속 수 하나하나)마다 풀어 쓰면, 이 식은 다음과 같아요.
\begin{equation}
r_j = d\sum_{i:(i\to j)\in E} r_i \frac{w_{ij}}{o_i} + \frac{d}{|V|}\sum_{i:\,o_i=0} r_i + \frac{1-d}{|V|}.
\label{eq:pagerank-node}
\end{equation}

\dissertationSubsubheading[sub:random-surfer]{무작위로 돌아다니는 사람으로 풀이하기}

식~\eqref{eq:pagerank-node}의 뜻은 확률로 풀어 읽을 수 있어요. 걸음마다, 나가는 화살표가 있는 점에 선 무작위로 돌아다니는 사람은 확률 $d$로 그 화살표 가운데 하나를 따라가요. 어느 화살표를 고를지는 그 무게에 비례해요. 아니면 확률 $1-d$로, 모든 점 가운데 고르게 뽑은 한 점으로 순간 이동(화살표를 무시하고 아무 점으로나 건너뛰기)해요(그림~\ref{fig:random-surfer}). 나가는 화살표가 없는 점에서는 언제나 고르게 뽑은 점으로 뛰어가요. 이것이 식~\eqref{eq:transition}의 매달린 점 규칙이에요. 그러면 정상 분포(오래 걸었을 때 각 점에 머무는 비율)는 오랫동안 걸었을 때 점마다 얼마나 자주 들르는지로 점에 순위를 매겨요. 감쇠 계수는 순위 웅덩이(밖으로 나가는 화살표가 없는 점 무리)가 몫을 모두 빨아들이지 못하게 막아요. $d<1$이고 순간 이동이 고르게 일어나면, $G_{\mathrm{PR}}$의 모든 칸이 양수(0보다 큰 수)예요. 그래서 $G_{\mathrm{PR}}$은 원시 행렬(몇 번 거듭 곱하면 모든 칸이 양수가 되는 행렬)이고, 그 페론 벡터(Perron vector, 모든 칸이 양수인 특별한 벡터)는 하나뿐이에요 \parencite{langville2006pagerank}. 그 논증은 부록~\ref{sec:pr-existence}에 있어요.

\begin{figure}[htbp]
\centering
\begin{tikzpicture}[
  font=\small,
  vertex/.style={draw, circle, minimum size=0.95cm},
  arrow/.style={-{Stealth[length=2.5mm]}, thick},
  teleport/.style={arrow, dashed}
]
\node[vertex] (a) at (0,0) {$A$};
\node[vertex] (b) at (3.2,1.5) {$B$};
\node[vertex] (c) at (3.2,-1.5) {$C$};
\node[vertex] (d) at (6.4,0) {$D$};
% Graph transitions: A->B, A->C, B->D, C->D (same topology as the worked example; D is dangling)
\draw[arrow] (a) -- node[above left, font=\footnotesize] {$2$} (b);
\draw[arrow] (a) -- node[below left, font=\footnotesize] {$1$} (c);
\draw[arrow] (b) -- node[below left=-1pt, font=\footnotesize] {$3$} (d);
\draw[arrow] (c) -- node[above left=-1pt, font=\footnotesize] {$1$} (d);
% Uniform teleportation, drawn from D only: one dashed arrow to every node, including D itself.
% D->A runs through the gap between B and C, and D->B / D->C bend outward, so no dashed arrow crosses a solid edge.
\draw[teleport] (d) -- (a);
\draw[teleport] (d) to[bend right=32] (b);
\draw[teleport] (d) to[bend left=32] (c);
\draw[teleport] (d) to[loop right, min distance=12mm] (d);
\end{tikzpicture}
\caption[{\ref{sec:pagerank-worked-example}절의 점 네 개짜리 그래프에서, 무작위로 돌아다니는 사람의 눈으로 페이지랭크를 풀어 본 그림이에요.}]{\ref{sec:pagerank-worked-example}절의 점 네 개짜리 그래프에서, 무작위로 돌아다니는 사람의 눈으로 페이지랭크를 풀어 본 그림이에요. 화살표에 붙은 수는 허락의 무게예요. 실선 화살표는 그래프를 따라 옮겨 가는 길이에요. 돌아다니는 사람은 모두 합쳐 확률 $d$로 그 가운데 하나를 따라가고, 어느 것을 고를지는 무게에 비례해요. 점선 화살표는 고른 순간 이동을 나타내요. 순간 이동은 확률 $(1-d)/|V|$로 모든 점(지금 있는 점도 포함)에 닿아요. 점선 부채꼴은 $D$에서만 그렸지만, 사실은 모든 점에 하나씩 있어요. $D$에는 나가는 화살표가 없어요. 그래서 식~\eqref{eq:transition}의 매달린 점 규칙에 따라, $D$에 있는 사람은 언제나 고르게 뽑은 점으로 뛰어가요. 각 점으로 갈 확률은 $1/|V|=0.25$예요. 감쇠 계수가 아니라 바로 이 규칙이, $D$에 닿은 몫을 걷기 안에 계속 남아 있게 해요.}
\label{fig:random-surfer}
\end{figure}

\dissertationSubsubheading[sub:power-iteration]{거듭제곱 반복과 계산 복잡도}

실제로는 보통 거듭제곱 반복(같은 계산을 되풀이하기)으로 $\mathbf{r}$을 구해요. 고른 벡터 $\mathbf{r}^{(0)}=\mathbf{1}/|V|$에서 시작해서 다음 계산을
\begin{equation}
\mathbf{r}^{(t+1)} = G_{\mathrm{PR}}\mathbf{r}^{(t)},
\label{eq:power-iteration}
\end{equation}
$\|\mathbf{r}^{(t+1)}-\mathbf{r}^{(t)}\|_1 < \varepsilon$이 될 때까지, 또는 반복 횟수 상한 $I_{\max}$에 이를 때까지 되풀이해요. 듬성듬성한 그래프에서는 반복 한 번에 $O(|E|)$번의 셈이 들어요 \parencite{leskovec2020mining}. 그러므로 화살표를 더 적게 만드는 방법은 반복 한 번에 드는 비용을 줄이고, 때로는 필요한 메모리도 줄여요. 송금 기록 전체 대신 마지막 허락(가장 최근 허락)의 스냅숏(어느 한 순간의 모습)을 쓰는 방법이 그 예예요.

이 연구는 기본으로 감쇠 계수를 흔히 쓰는 값인 $d=0.85$로 정해요 \parencite{page1999pagerank}. 그리고 이 연구에서만 정한 풀이 프로그램(solver) 설정 두 가지를 써요. 수렴(값이 더는 거의 바뀌지 않게 됨)의 허용 오차 $\varepsilon=10^{-8}$과, 반복 횟수 상한 $I_{\max}=300$이에요.\footnote{허용 오차와 반복 횟수 상한은 실행 시간을 재현할 수 있게(다른 사람이 똑같이 다시 해 볼 수 있게) 하려고 골랐어요. 이 둘은 고전적인 페이지랭크 정의의 한 부분으로 내놓은 것이 아니에요.} 강건성 점검(조건을 바꿔도 결과가 버티는지 보기)을 다루는 절에서는 감쇠 계수를 $d\in\{0.75,0.85,0.95\}$ 안에서 바꿔 봐요.

\dissertationSubheading[sec:reputation-systems]{평판 시스템이 이어져 온 흐름}

페이지랭크는 스펙트럼 방법(행렬의 고유벡터로 점수를 구하는 방법)과 퍼뜨림에 바탕을 둔 방법으로 평판을 매기는, 더 넓은 방법 가족에 속해요. Kleinberg의 HITS 알고리즘은 페이지마다 점수 두 개를 줘요. 하나는 허브(좋은 페이지를 많이 가리키는 페이지)로서의 점수이고, 다른 하나는 권위(좋은 페이지들이 많이 가리키는 페이지)로서의 점수예요 \parencite{kleinberg1999hits}. EigenTrust는 이 생각을 P2P(컴퓨터끼리 직접 주고받는) 네트워크로 넓히고, 이웃끼리 주는 믿음을 정규화(합이 1이 되게 맞추기)해요. 그래서 나쁜 뜻을 가진 점이 제 점수를 마음대로 올릴 수 없어요 \parencite{kamvar2003eigentrust}. TrustRank는 웹 스팸(검색 순위를 속이려는 가짜 페이지)을 누르려고, 믿을 수 있는 점들로 이루어진 씨앗 집합(출발점으로 골라 둔 점들)에서 출발해 그 믿음을 바깥으로 퍼뜨려요 \parencite{gyongyi2004trustrank}. 이 모든 방법에서 평판은 관계에서 생겨요. TrustRank처럼 꼼꼼히 골라 둔 씨앗 집합에서만 믿음이 흘러나오면, 새로 온 점은 이미 평판을 가진 점들의 보증(믿고 맡김) 없이는 점수를 올릴 수 없어요. 엔도스랭크는 고른 순간 이동을 써요(\ref{sub:endorserank-graph}절). 이때는 그 보호가 더 약해요. \textcite{gyongyi2005alliances}에서는 목표 점의 페이지랭크를 가장 많이 올리는 링크 농장(점수를 올리려고 서로 링크를 거는 가짜 페이지 무리)과 농장 동맹(여러 농장이 손잡은 것)을 끌어냈어요. \textcite{cheng2005sybilproof}에서는 페이지랭크 같은 대칭 평판 함수(모든 점을 똑같이 다루는 함수)를 시빌로 조작할 수 있음을 보였어요. 반면에 믿을 수 있는 점에서 나오는 흐름이나 경로에 바탕을 둔 어떤 비대칭 함수들은 시빌을 견뎌 낸다는 것도 보였어요. \ref{sec:sybil-cost-model}절의 시빌 비용 모형은 허락 그래프 위에서 이런 농장을 만드는 데 드는 값을 매겨요.

온체인에서는 시빌 공격이 여전히 중심이 되는 걱정거리예요 \parencite{bordeianu2026sybil}. 새 지갑은 거의 돈을 들이지 않고 만들 수 있어요. 하지만 서로 다른 거래 상대들에게서 허락이나 송금을 받는 지갑은 만드는 데 비용이 더 들어요. 이 연구에서 쓰는 시빌 보정 안정성 대리 지표(대신 재는 값)는 들어오는 거래 상대 비율, 활동 기간, 활동한 달 수예요. 이 지표들은 이런 주고받음의 다양함을 어느 정도 담아요. 하지만 나쁜 뜻을 가진 사람이 이 지표들을 조작하는 것을 막지는 못해요.

디파이 대출 자료에 맞춘 기계 학습 신용 모형은 라벨이 붙은 신용 결과에서 배워요 \parencite{kotb2023credit}. 그리고 맞춘 모형 여러 개를 묶은 앙상블(여러 모형을 합친 것)은 정해진 규칙보다 감사하기(꼼꼼히 따져 확인하기)가 더 어려워요. 엔도스랭크와 AWP는 정해진 퍼뜨림 규칙을 지켜요. 누구나 공개된 로그로 이 규칙을 다시 계산할 수 있어요.

\dissertationSubheading[sec:allowance-endorsement]{보증으로 보는 ERC-20 허락}

ERC-20의 \texttt{approve} 함수 \parencite{vogelsteller2015erc20}와 EIP-2612의 \texttt{permit} \parencite{lundfall2020permit}을 써서, 사용자는 스펜더가 자기 대신 얼마까지 옮겨도 되는지 분명하게 밝혀요. 엔도스랭크는 주인이 스펜더에게 준 마지막 허락을 보증 하나로 읽어요. 이 연구의 토큰 다섯 개(\ref{sec:amount-revision}절) 각각에서, 그 쌍의 기간 안 마지막 허락이 항 하나를 보태요. 그 항은 시간 무게(오래될수록 가벼워지는 무게)에, 그 허락의 달러 금액에 매긴 상한이 있는 금액 무게(금액이 아무리 커도 1을 넘지 않는 무게)를 곱한 값이에요. AWP가 송금에 무게를 매기는 것과 같아요. 마지막 값이 0인 쌍은 빼요. 이것은 허락 취소이거나, \texttt{transferFrom}이 \texttt{Approval}을 남기는 토큰에서는 다 써 버린 허락이에요. 그리고 이 항들을 토큰에 걸쳐 더해요(\ref{sub:endorserank-graph}절의 식~\eqref{eq:approve-weight}). 그렇게 만든 그래프 $G_{\text{approve}}$는 같은 계약들과 같은 기간으로 만든 송금 그래프보다 더 듬성듬성해요. 이 그래프에 페이지랭크를 돌리면 엔도스랭크 점수 $\mathbf{s}_{\mathrm{ER}}$이 나와요.

허락을 보증으로 보는 것은 일부러 한 선택이고, 꼭 들어맞지는 않는 선택이에요. 사용자가 라우터를 허락하는 것은 편해서이지, 꼭 그 라우터가 믿을 만하다고 남들 앞에서 말하려는 것은 아니에요. Uniswap의 Permit2 같은 허락 관리자(여러 앱의 허락을 한데 모아 다루는 계약)도 있어요. 주인들은 이 관리자를 한 번 허락해 두고, 앱들이 이 관리자를 거쳐 토큰을 쓰게 해요. 그래서 이 관리자는 원래 그 앱들에게 가야 할 보증을 대신 모아요(\ref{sub:endorserank-graph}절). 또 주인 쪽에도 계약 주소와 외부 소유 계정(EOA, 개인 키로 움직이는 계정)이 함께 있어요. 그래도 허락은 받은 송금보다 찬성의 뜻을 더 분명하게 보여 줘요. 받은 송금은 시장 조성(늘 사고팔 물량을 내놓는 일), 자전 거래 \parencite{cong2023washtrading}, 에어드롭 파밍 \parencite{luo2025airdrop}, 또는 보증할 뜻이 전혀 없는 경로 이어 주기에서 생길 수도 있기 때문이에요.

허락을 주소 사이의 관계로 읽은 연구는 전에도 있었어요. Victor의 자기 허락 어림 방법(휴리스틱)은 ERC-20 허락을 써서, 한 주체의 것이라고 보는 주소들을 한데 묶어요 \parencite{victor2020clustering}. 이렇게 읽으면 주소들을 이을 뿐, 순위는 하나도 매기지 않아요. 이것은 조심하라는 신호이기도 해요. 한 주체가 가진 두 주소 사이의 허락은 자기가 자기를 보증하는 것이기 때문이에요. \ref{sec:sybil-cost-model}절의 시빌 비용 모형은 이런 화살표의 값을 가스 비용만으로 매겨요.

\dissertationSubheading[sub:non-claims]{이 연구가 주장하지 않는 것}

거래 결과는 이 논문에서 기간 밖(점수를 매긴 기간이 끝난 뒤)에서만 쓰여요. 먼저 W0(첫째 기간)의 GMX~V2 계약 트레이더(GMX V2에서 거래하는 계약)에게 해 본 탐색적(미리 정하지 않고 살펴본) 실행에 쓰여요(\ref{sub:trader-results}절). 다음으로 등록 기간 W1(계획을 먼저 등록해 둔 기간)의 청산 없이 닫은 비율(강제로 끝나지 않고 거래를 닫은 비율)로 쓰여요(\ref{sub:fresh-holdout}절). 그리고 두 종류의 화살표를 견주는 비교에 쓰여요(\ref{sub:neutral-label-results}절). 그 비교에는 두 종류의 화살표 어디에서도 만들지 않은 결과가 필요해요. 구성 타당도(점수가 재려는 것을 정말 재는지)에서는, 점수가 자기가 재도록 만들어진 구성개념(점수가 재려고 하는 생각)의 점검과 잘 맞기를 기대해요 \parencite{cronbach1955construct,messick1995validity}. 그래서 화살표 종류마다 자기 점검에서 앞서는 것은, 어느 정도 만든 방식 때문에 저절로 생겨요. GMX~V2 트레이더의 결과에서는 어느 점수도 그렇게 앞서서 출발하지 못해요. 그리고 거래 상대 위험을 다루는 평판에서 따져야 할 결과는 강제 청산이에요. 이 논문의 어떤 점수도 거래 이익, 손실을 피하는 것, 청산을 예측하는 값으로 내놓지 않아요. 신용도(돈을 빌려줘도 될 만한 정도)나 빚을 못 갚을 위험을 재는 값으로도 내놓지 않아요. 또 어떤 점수도 계약 코드의 보안 감사(코드에 위험한 곳이 없는지 꼼꼼히 살피기)가 아니에요. 계약은 주인들 사이에서 평판이 어떻든 나쁜 함수를 품고 있을 수 있어요 \parencite{lin2024crpwarner}. 그리고 점수는 코드도, 업그레이드 권한(계약을 고쳐 바꿀 수 있는 권한)도 읽지 않아요.

\dissertationSubheading[sec:rank-correlation-theory]{순위 상관 이론}

평판 점수와 구성개념 점검의 대리 지표는 둘 다 순서를 매길 수 있는 값이에요. 그리고 어느 쪽에도 둘로 나눌 자연스러운 경계가 없어요. 이 연구는 AWP의 검증 설계 \parencite{do2023awp}와 고전 심리 측정학(사람의 마음과 능력을 수로 재는 학문) \parencite{cronbach1955construct}을 따라, 두 가지 순위 일치를 재요. 하나는 단조 일치(한쪽이 커질 때 다른 쪽도 커지는지)이고, 다른 하나는 쌍 일치(둘씩 짝지었을 때 순서가 같은지)예요. 부트스트랩(다시 뽑기) 신뢰구간(믿을 만한 범위)은 주 분석의 모든 계수가 표본에 따라 얼마나 흔들리는지(표본 불확실성)를 수로 나타내요 \parencite{efron1979bootstrap}.

스피어먼의 순위 상관은 순위들끼리 구한 피어슨 상관(두 값이 함께 움직이는 정도)이에요 \parencite{spearman1904proof}. 동점(값이 같은 것)이 없는 $n$개 항목의 짝지은 순위 $(R_i,S_i)$에서는, 이 값이 다음처럼 간단해져요.
\begin{equation}
\rho = 1 - \frac{6\sum_{i=1}^n (R_i-S_i)^2}{n(n^2-1)}.
\label{eq:spearman}
\end{equation}
동점이 있으면 이 동점 없는 식은 맞지 않아요. 그때 스피어먼의 로($\rho$)는 평균 순위(가운데 순위)의 피어슨 상관이에요. 평균 순위에서는 동점인 항목들이 자기들이 차지한 자리들의 평균을 함께 나눠 가져요. 이 논문에서 보고하는 $\rho$는 모두 이렇게 계산했어요. 두 식 모두에서 $\rho$는 단조 연관(한쪽이 커질 때 다른 쪽도 커지는 관계)을 재요.

켄달의 타우($\tau$, 두 줄 세우기가 얼마나 같은 순서인지 나타내는 수)는 순서가 맞는 쌍과 어긋나는 쌍의 수를 바탕으로 해요 \parencite{kendall1938rank}. $i<j$인 쌍 $(i,j)$는 모두 $n_0=n(n-1)/2$개예요. 그 가운데 $n_c$개는 순서가 맞는 쌍이에요(두 순위가 같은 쪽으로 줄 세워요). $n_d$개는 어긋나는 쌍이에요(두 순위가 반대로 줄 세워요). 어느 한 순위에서라도 동점인 쌍은 둘 다 아니에요. $n_1=\sum_k t_k(t_k-1)/2$라고 해요. 여기서 $t_k$는 첫째 순위에서 $k$번째 동점 무리의 크기예요. $n_2$도 둘째 순위에 대해 똑같이 정해요. 그러면 $\tau$-b 변형은 다음과 같아요.
\begin{equation}
\tau_b = \frac{n_c - n_d}{\sqrt{(n_0-n_1)(n_0-n_2)}},
\label{eq:kendall}
\end{equation}
그리고 동점 보정을 빼면 이 식은 $\tau_a=(n_c-n_d)/n_0$ 꼴이 돼요. 이 논문에서 보고하는 $\tau$는 모두 $\tau_b$이고, 손대지 않은 원래 점수와 대리 지표로 계산했어요. 이 차이는 동점이 아주 많은 점수에서 중요해요. 정수로 센 값인 원시 차수(화살표를 그냥 센 수)가 그런 점수예요. 매칭 코호트(모든 점수가 똑같이 점수를 매기는 계약 묶음)에서 엔도스랭크는 동점을 거의 남기지 않아요(\ref{sec:rank-divergence}절). 둘씩 짝지은 순서가 얼마나 일치하는지 볼 때는, 보통 켄달의 $\tau$가 더 해석하기 쉬운 계수예요. 그래서 \ref{ch:results}장은 이 계수를 지표 가족(비슷한 지표 묶음) 수준의 주된 요약으로 써요. 그 값은 같은 코호트에서 다른 방법들의 값과 견주어요. 출처를 밝히지 않은 절대 기준값은 쓰지 않아요.

\dissertationSubheading[sec:theoretical-framework]{이론의 틀}

이 연구에서는 먼저 한 층만 걷는 두 걷기, 곧 엔도스랭크와 AWP를 평판 구조(H1), 구성개념 일치도(H2), 계산 효율(H3)에서 견주어요. 그리고 기간 밖에서 모든 점수를 원시 차수와 견주어요(RQ3). 그다음에는 두 걷기의 층(두 종류의 화살표)을 결합하는 걷기를 시험해요(RQ4와 RQ5). 앞부분은 세 가지 이론이 틀을 잡아 줘요. 하이퍼링크(웹 페이지끼리 잇는 링크)--인용 퍼뜨림, 구성개념(분야) 일치도(점수가 재려는 것과 얼마나 잘 맞는지), 계산 복잡도예요. 뒷부분에는 두 가지 논리가 더 관련돼요.

\begin{itemize}
\item 하이퍼링크--인용의 뜻: 원래 페이지랭크 알고리즘에서는 들어오는 하이퍼링크를 인용 하나로 세어요 \parencite{page1999pagerank}. \textcite{do2023hodgerank}에서는 이 생각을 이더리움 송금에 썼어요. AWP는 송금 화살표마다 그 송금들의 나이와 금액으로 무게를 달아요 \parencite{do2023awp}. 나이에는 로지스틱 감쇠(오래된 일일수록 가볍게 세는 무게)를 쓰고, 금액에는 상한이 있는 값 변환을 써요(\ref{sub:awp-graph}절). 엔도스랭크는 인용에 바탕을 둔 퍼뜨림을 ERC-20 허락하기 화살표와 서명 허락 화살표로 넓혀요. 허락은 쓸 권한을 온체인에서 분명하게 맡기는 일이에요. 이 허락은 다 쓰이거나, 거두어지거나, 다른 허락으로 바뀔 때까지 효력이 있어요 \parencite{vogelsteller2015erc20}. 이와 달리 송금은 믿으려는 뜻 없이 그저 거래하거나 길을 이어 준(라우팅) 기록일 수도 있어요. 그래프에 바탕을 둔 위험 모형은 거래 연결망을 따라 위험을 퍼뜨려서, 라벨이 없는 계정에도 등급을 매겨요 \parencite{lin2025riskprop}. 여기서 다루는 질문은, 방향이 있는 권한 위임 그래프가 송금 위의 이런 퍼뜨림에 무언가를 보태는지예요.

\item 분야 일치도의 틀: 측정 이론(무엇을 어떻게 재는지 다루는 이론)에서는 구성개념 일치도로, 한 지표가 따라가려는 구성개념을 그 지표가 얼마나 가깝게 따라가는지 물어요 \parencite{cronbach1955construct}. \textcite{messick1995validity}에서는 이 생각을 더 넓혔어요. 점수가 무엇을 뜻하는지, 그리고 그 점수를 쓰면 무엇이 따르는지를 하나로 묶어 판단하게 했어요. 그래프가 앞뒤가 맞는지(일관성)는, 같은 기간의 송금 통계와 허락 통계에 견준 순위 상관으로 점검해요(\ref{sec:rank-correlation-theory}절). 기간 밖 점검은 시간 홀드아웃(나중 기간을 떼어 두고 맞혀 보는 시험)이에요. 스펜더에서는 새 허락과 새 송금자(새로 송금을 보내온 주소)를 봐요. 그리고 7월부터 9월까지의 등록 기간 W1에서는, 트레이더의 청산 없이 닫은 비율을 봐요.

\item 계산 복잡도: 페이지랭크의 반복 한 번에는 $O(|E|)$만큼의 계산이 들어요(\ref{sub:power-iteration}절). 그래서 더 듬성듬성한 허락 그래프에서는 반복 한 번에 드는 시간과 메모리가 더 적어요. 이 논리가 H3의 바탕이에요. 반복 횟수는 이와 따로인 요인이에요. 감쇠 계수는 반복이 수렴하는 빠르기의 한계를 정해요 \parencite{langville2006pagerank}. 하지만 반복 횟수는 그래프와 재시작 벡터(어디서 다시 출발하는지 정하는 값들)에도 달려 있어요. 두 방법은 같은 풀이 프로그램을 써요. 그래서 실행 시간의 차이는 그래프와 재시작 벡터에서 나오고, 알고리즘의 기여(이 연구가 새로 보탠 것)로 치지 않아요(부록~\ref{sec:pr-sparsity}).

\item 기준선이 되는 국소 개수(바로 이어진 화살표만 센 수): 이미 거래 상대가 많은 계약이 새 상대도 가장 많이 얻는다면(\ref{sec:graph-reputation}절), 원시 차수는 앞으로 새로 들어올 상대를 잘 예측해요. 페이지랭크는 원시 차수와 세 가지가 달라요. 보증해 주는 쪽 자신의 점수가 화살표를 따라 퍼지고, 화살표에 나이와 금액으로 무게를 달고, 걷기가 자기만의 규칙으로 다시 출발해요. 차수와 대비해 보면 이 세 가지를 한꺼번에 시험하게 돼요. 그래서 RQ3에서는 각 페이지랭크를, 동결일(점수를 고정하는 날)에 잰 같은 종류 화살표의 원시 차수와 견주어요. 그리고 차이가 나더라도 그것을 퍼뜨림 하나만의 덕으로 돌리지 않아요.

\item 층 결합: 다중 페이지랭크(multiplex PageRank, 여러 층을 함께 보는 페이지랭크) \parencite{halu2013multiplex}와 다층 순위 매기기 \parencite{dedomenico2015ranking}는 여러 관계를 한 점 집합 위의 층들로 다뤄요. C-PR(결합 페이지랭크, 두 그래프를 함께 걷는 점수)처럼 점마다 층들을 섞으면, 빽빽한 정도가 아주 다른 층들에서는 풀리지 않은 질문이 남아요. 듬성듬성한 층에 화살표가 없는 점은 빽빽한 층만 따라가요. 그래서 듬성듬성한 층의 신호가 살아남을 수도 있고, 사라질 수도 있어요. 이 때문에 RQ4와 RQ5에서는 두 가지를 물어요. 섞은 것이 두 층을 모두 이기는지, 그리고 송금 층만 걸을 때보다 나은지예요. 두 층은 모두 AWP의 상한이 있는 금액 무게 없이 달러 금액으로 화살표에 무게를 달고, 고르게 다시 출발하기(균등 재시작)로 걸어요. 그래서 어느 걷기도 AWP나 엔도스랭크가 아니에요. C-PR은 AWP와 따로 견주어요(\ref{sub:cpr-contrasts}절). 거기서는 덮는 범위(점수로 순서를 매길 수 있는 점이 얼마나 되는지)가 중요해요. 송금 걷기는 송금 화살표가 없는 점들을 모두 한 동점 덩어리로 남겨요. 그런데 허락 화살표 위의 걷기는, 엔도스랭크를 포함해 어느 것이든 그 점들에 순서를 매겨요. 그래서 결합은 두 층이 모두 닿는 곳에서는 아무것도 보태지 않더라도, 송금 걷기보다 덮는 범위를 넓힐 수 있어요.
\end{itemize}

그림~\ref{fig:conceptual-model}의 개념 모형(생각을 그림으로 정리한 틀)은 이 논문의 점수들을 H1--H3, RQ1--RQ5와 이어요. 한 층 걷기 두 가지는 H1--H3, RQ1과 RQ2, 그리고 원시 차수와 견주는 홀드아웃(RQ3)에 쓰여요. 오른쪽의 점선 상자는 이 논문의 방법인 결합 걷기 C-PR이에요. C-PR은 구성개념 점검과 홀드아웃에서 입력을 받고, 두 구성개념이 합쳐지는지 물어요(RQ4와 RQ5).

\begin{figure}[htbp]
\centering
% !TEX root = ../main.tex
\begin{tikzpicture}[
  font=\small,
  box/.style={draw, rounded corners, align=center, minimum width=10.5cm, minimum height=1.0cm},
  smallbox/.style={draw, rounded corners, align=center, minimum width=3.4cm, minimum height=1.05cm},
  refbox/.style={draw, dashed, rounded corners, align=center, minimum width=3.4cm, minimum height=0.95cm},
  arrow/.style={-{Stealth[length=2.5mm]}, thick}
]
\node[box] (center) at (0,3.2) {계약 코호트 하나에 매긴 허락 점수와 송금 점수};

\node[smallbox] (rep) at (-5,1.2) {평판\\구조};
\node[smallbox] (risk) at (0,1.2) {구성개념\\점검};
\node[smallbox] (eff) at (5,1.2) {계산\\비용};

\node[smallbox] (h1) at (-5,-1.0) {H1 / RQ1\\엔도스랭크--AWP\\갈라짐};
\node[smallbox] (h2) at (0,-1.0) {H2 / RQ2\\각 구성개념과의\\일치};
\node[smallbox] (h3) at (5,-1.0) {H3 / RQ3\\각 점수의\\비용};

\node[refbox] (hold) at (0,-3.0) {RQ4: 시간 홀드아웃\\점수 대 원시 차수};
\node[refbox] (rq5) at (5,-3.0) {RQ5: 결합 연산자 C-PR\\두 화살표 종류, 두 라벨 모두};

\draw[arrow] (center.south -| rep.north) -- (rep.north);
\draw[arrow] (center.south) -- (risk.north);
\draw[arrow] (center.south -| eff.north) -- (eff.north);

\draw[arrow] (rep.south) -- (h1.north);
\draw[arrow] (risk.south) -- (h2.north);
\draw[arrow] (eff.south) -- (h3.north);

\draw[arrow, dashed] (h2.south) -- (hold.north);
\draw[arrow, dashed] (h2.south east) -- (rq5.north west);
\draw[arrow, dashed] (hold.east) -- (rq5.west);
\end{tikzpicture}

\caption[{점수들을 H1--H3, RQ1--RQ5와 이어 주는 개념 모형이에요.}]{계약 코호트 하나의 허락 점수와 송금 점수를 평판 구조와 구성개념 점검(H1, H2/RQ1), 계산 비용(H3/RQ2), 시간 홀드아웃에 이어 주는 개념 모형이에요. 시간 홀드아웃에서는 각 페이지랭크를, 그 페이지랭크가 매끄럽게 다듬는(평활하는) 원시 차수와 견주어요(RQ3). 오른쪽의 점선 상자는 이 논문의 방법인 결합 걷기 C-PR이에요. C-PR은 규칙 A와 규칙 B로 시험해요(RQ4, RQ5).}
\label{fig:conceptual-model}
\end{figure}

H1에는 엔도스랭크와 AWP만 들어가요. H2와 H3은 모든 점수에 관한 것이에요. RQ3은 스펜더 코호트에서 시간으로 나눠 하는 시험이에요. 동결 때의 원시 차수들을 기준선으로 삼고, 거래 결과는 줄 세우지 않아요. RQ4와 RQ5는 규칙 A와 규칙 B로 C-PR을 시험해요.

\dissertationSubheading[sub:terminology]{용어}

정식 구성개념은 \ref{ch:introduction}장에서 정의해요. 아래 용어들은 결과를 풀이할 때 써요.

\begin{itemize}
\item 한 층 방법: 엔도스랭크(허락 그래프)와 AWP(송금 그래프)예요. 둘 다 한 종류의 화살표 위에서만 걷는 페이지랭크 점수예요. \ref{sub:coupled-graph}절의 결합 걷기들과 원시 차수들도 이 둘과 나란히 평가해요.
\item 주 그래프 구성개념: 한 방법이 순위를 매기도록 \emph{설계된} 그래프예요. 예를 들어 엔도스랭크에서는 $G_{\text{approve}}$예요. 이것은 점수를 매긴 뒤에 하는 점검과는 달라요.
\item 같은 기간 점검: 점수와 같은 기간에서 가져온 국소 송금 통계나 국소 허락 통계예요 \parencite{cronbach1955construct}.
\item 시간 홀드아웃: 첫째 기간 W0에서는 점수를 $t_1$(2026년 3월 31일)에 동결해요. 그리고 새 주인--스펜더 허락 쌍과 새 송금자를 $t_2$(2026년 6월 30일)까지 세어요. 등록한 둘째 기간 W1은 2026년 6월 30일에 점수를 동결하고, 2026년 7월부터 9월까지 라벨을 세어요(\ref{sub:fresh-holdout}절).
\item 매칭 계약 코호트, W0과 W1의 스펜더 홀드아웃 코호트, 트레이더 코호트, 매칭 코호트의 부분 표본: \ref{sub:cohorts}절의 주소 집합들이에요.
\end{itemize}

\dissertationSubheading[sec:domain-alignment-metrics]{분야 일치도 지표}

구성개념 일치도는 점수와 점검 순위 사이의 스피어먼의 $\rho$와 켄달의 $\tau$로 재요(\ref{sec:rank-correlation-theory}절). 이때 같은 기간의 송금 점검과 허락 점검, 그리고 홀드아웃 라벨을 써요. 실제로 어떻게 계산하는지(조작적 정의)는 \ref{ch:methodology}장에서 알려 줘요.

\dissertationSubheading[sec:pagerank-worked-example]{풀어 본 예: 점이 네 개인 보증 그래프에서의 페이지랭크}

주소 네 개 $\{A,B,C,D\}$가 아래의 마지막 허락 화살표로 이어져 있다고 해요.
\[
A\to B\ (2),\quad A\to C\ (1),\quad B\to D\ (3),\quad C\to D\ (1),
\]
감쇠 계수는 $d=0.85$예요. 나가는 무게의 합은 $o_A=3$, $o_B=3$, $o_C=1$, $o_D=0$이에요(마지막 것은 매달린 점이에요). 화살표 $B\to D$와 $C\to D$는 각자 출발하는 점에서 나가는 하나뿐인 화살표예요. 그래서 각 화살표는 반드시(확률 1로) 따라가게 되고, 그 무게는 점수에 아무 영향도 주지 않아요. 순간 이동은 첫 단계를 포함해 단계마다 모든 점에 $(1-d)/4=0.0375$를 더해요. 고른 벡터 $\mathbf{r}^{(0)}=\mathbf{1}/4$에서 거듭제곱 반복을 한 단계 하면, $(A,B,C,D)$에 대해 $\mathbf{r}^{(1)}=(0.0906,\,0.2323,\,0.1615,\,0.5156)$을 얻어요. 벌써 몫이 $D$ 쪽으로 기울어요. $B$와 $C$가 둘 다 $D$를 보증하기 때문이에요. 한편 $A$는 자기 몫을 $B$와 $C$에 나눠 줘요. 반복은 정상 벡터 $\mathbf{r}=(0.1375,\,0.2154,\,0.1765,\,0.4706)$으로 수렴해요. 그래서 $D$가 맨 앞이고, $B$와 $C$가 그 뒤를 따르고, $A$가 맨 뒤예요. $A$를 보증하는 주소는 하나도 없어요. 순간 이동을 빼면, 이 주소의 점수는 모두 매달린 점 규칙에서 나와요. 이 규칙은 $D$의 몫을 네 점에 고르게 나눠 줘요($d\,r_D/4=0.1000$). 이 예는, 한 점이 자기도 보증을 받는 점들에게서 보증을 \emph{받아서} 점수를 얻는다는 것을 보여 줘요. 이것이 바로 페이지랭크의 인용 논리예요 \parencite{page1999pagerank,langville2006pagerank}.

같은 주소들이 허락 대신 송금으로 이어져 있다고 해 봐요. 그러면 송금은 보통 다른 쌍들을 잇거나, 같은 쌍이라도 쌍마다 합계가 다를 거예요. 그러면 이동 행렬 $P$가 달라지고, 같은 점들에 다른 순위가 매겨질 수도 있어요. 가설 H1은 이 상황을 작게 줄여 놓은 것이에요. 엔도스랭크와 AWP가 꼭 같은 순위를 낼 필요는 없어요. 두 방법은 다시 출발하는 방식도 달라요. 그래서 H1은 화살표와 다시 출발을 함께 견주어요. \ref{sub:endorserank-restarts}절에서는 AWP의 다시 출발 방식으로 허락 그래프를 걸어 봐요.

\dissertationSubheading[sec:defi-risk-taxonomy]{디파이 위험의 분류와 평판이 들어갈 자리}

\textcite{werner2022sok}에서는 디파이 서비스의 기술적 보안과 경제적 보안을 나누었어요. 그리고 \textcite{zhou2023sok}에서는 실제로 일어난 공격과 사고를 하나의 기준 틀 안에 정리했어요. 계약 코드의 버그(프로그램 오류)는 스마트 계약 위험이에요. 이것은 디파이의 약점(취약점)이 생기는 여러 수준 가운데 하나예요 \parencite{li2022defisecurity}. 오래되어 맞지 않는 가격이나, 누군가가 마음대로 다루는 가격 피드(블록체인에 올라오는 가격 정보)는 오라클 위험이에요 \parencite{angeris2020oracles}. 시장 위험에는 가격이 크게 출렁이는 것(변동성)과 청산이 줄줄이 이어지는 것이 들어가요 \parencite{perez2021liquidations}. 스테이블코인(값이 거의 고정되도록 만든 토큰)과 자금 시장이 받는 압박도 여기에 들어가요 \parencite{klagesmundt2020stablecoins}. 2020년 3월의 시장 폭락은 이 두 층이 서로 어떻게 영향을 주고받을 수 있는지 보여 주었어요 \parencite{gudgeon2020crisis}. 가진 토큰 수만큼 표를 주는 투표는 거버넌스 위험(서비스의 규칙을 정하는 방식에서 오는 위험)에 들어요. Compound, Uniswap, ENS의 온체인 거버넌스에서는 적은 수의 주소가 투표권의 대부분을 쥐고 있어요 \parencite{fritsch2024voting}. 우리는 더 좁은 층인 \emph{거래 상대} 위험, 곧 \emph{행동} 위험을 살펴봐요. 곧, 한 주소의 지난 행동으로 보아 그 주소가 안전한 거래 상대인지를 봐요. 레버리지(빌린 돈을 보태 크게 거는 것)를 책임감 있게 다루는 트레이더로서, 또는 많은 주인이 믿고 자기 토큰을 쓰게 해 주는 계약으로서 안전한지예요.

온체인 평판 점수는 행동 위험 층의 일부만 다뤄요. 그리고 담보, 감사, 오라클 설계를 대신하지 못해요. 평판 점수의 목적은 한 시스템의 관계 구조(누가 누구에게 허락하는지, 누가 누구에게 돈을 내는지)를 순서를 나타내는 값 하나로 줄이는 것이에요. 다른 사람들은 이 값을 공개된 자료로 확인하고, 똑같이 다시 만들어 볼 수 있어요. 행동 위험 안에서 엔도스랭크는 허락을 보고, AWP는 돈의 흐름을 봐요. 둘 다 취약점 공격(익스플로잇)이나 거버넌스 공격을 예측하려고 하지는 않아요.

마지막으로, 가짜 활동 때문에 행동 위험을 재기가 더 어려워요. 자전 거래는 탈중앙 거래소(중간 관리자 없이 계약으로 움직이는 거래소)의 거래량을 부풀려요 \parencite{victor2021washtrading}. 규제를 받지 않는 중앙 거래소에서는 더 크게 부풀려요 \parencite{cong2023washtrading}. 앞질러 거래하기(프런트러닝)와 이와 비슷한 추출 가능한 가치(거래 순서를 이용해 빼내는 이익)는, 관계가 아니라 순서 전략을 따르는 송금을 만들어요 \parencite{daian2020flashboys}. \textcite{qin2022bev}에서는 이더리움에서 그런 가치가 얼마나 빼내어지는지 어림했어요. Uniswap의 사기 토큰도 마찬가지로, 오래가는 관계를 만들지 않고 거래를 늘려요 \parencite{xia2021scamtokens}. 송금 그래프는 가짜 거래량에 특히 약해요. 이와 달리 허락 그래프는 허락 파밍(가짜 허락을 잔뜩 만들어 내기)에 약해요. 시간으로 나누기와 시빌 탐지는 이 문제를 일부만 풀어요.

\dissertationSubheading[sec:related-measurement]{관련된 측정 연구의 전통}

블록체인 밖에서는 전자 상거래 \parencite{resnick2002trust}, P2P 네트워크 \parencite{kamvar2003eigentrust}, 웹 스팸 연구 \parencite{gyongyi2004trustrank}에서 평판 시스템을 연구해 왔어요. 이 연구들에서 얻은 깨달음 네 가지가 여기에도 들어맞아요. (i) 평판은 관계가 가진 성질이에요. (ii) 신분을 쉽게 만들 수 있으면 시빌 공격이 가능해져요 \parencite{douceur2002sybil}. (iii) 씨앗 정하기와 정규화는 부풀리기의 효과를 줄여요. (iv) 평가에는 바깥 기준(벤치마크)이 필요해요. 안에서 앞뒤가 맞는 것만으로는 모자라기 때문이에요 \parencite{cronbach1955construct}. 엔도스랭크는 (i)을 바로 받아들이고, (ii)를 위협으로 다뤄요. (iii) 가운데 정규화는 받아들여요. 주인마다 자기 보증을 자기가 허락한 스펜더들에게 나눠 주기 때문이에요. 하지만 씨앗 정하기는 받아들이지 않아요. 고르게 다시 출발하기 때문이에요(\ref{sec:reputation-systems}절). (iv)는 일부만 따라요. 스피어먼의 $\rho$ \parencite{spearman1904proof}와 켄달의 $\tau$ \parencite{kendall1938rank}로 재는 같은 기간 점검은 앞뒤가 맞는지를 시험해요. 한편 시간 홀드아웃은 기간 밖 예측을 시험해요. 스펜더에서는 같은 구성개념으로 시험하고, 트레이더에서는 청산으로 시험해요.

개인화 페이지랭크와 주제 민감 페이지랭크 \parencite{haveliwala2003topic,jeh2003personalized}는 순간 이동 벡터를 골라 둔 점 무리 쪽으로 치우치게 해요. 트러스트랭크(TrustRank)도 이 방법으로 작은 씨앗 집합을 전체 순위로 바꿔요 \parencite{gyongyi2004trustrank}. 보내는 활동에 비례해서 다시 출발하는 AWP도 또 하나의 예예요 \parencite{do2023awp}. 다중 페이지랭크에서는 한 층에서 점이 얼마나 중심에 있는지(중심성)가 다른 층의 걷기를 바꿔요 \parencite{halu2013multiplex}. 그리고 \textcite{dedomenico2015ranking}의 다층 틀에서는 걷는 사람 한 명이, 서로 이어진 네트워크의 층 안에서도 움직이고 층 사이로도 움직여요. 두 방법 모두 층을 한 점 집합 위의 서로 다른 관계로 다뤄요. 이것은 같은 주소들 사이의 허락과 송금에 딱 맞아요. C-PR은 이 생각을 가장 단순한 꼴로 가져와요. 곧, 두 이동 행렬을 무게 하나로 점마다 섞어 더한 것이에요. 그리고 S-PR(보증 씨앗 페이지랭크)은 개인화 순간 이동이라는 생각을 빌려 와요. S-PR은 허락 점수를 순간 이동 벡터로 써요(\ref{sub:coupled-graph}절).

\dissertationSubheading[sec:account-model]{계정, 가스, 그리고 로그로 재는 까닭}

이더리움에는 두 종류의 계정이 있어요. 외부 소유 계정(EOA)과 계약 계정이에요. EOA는 개인 키(주인만 아는 비밀 열쇠)가 다루고, 계약 계정은 자기가 저장한 코드가 다스려요 \parencite{antonopoulos2018mastering}. 이 논문은 계약 계정에만 점수를 매겨요(\ref{sec:data-sources}절). 상태를 바꾸는 거래는 모두 가스를 써요. 로그는 계약의 \texttt{LOG} 명령에서 나와요. 이 명령은 계약의 코드가 이벤트를 내보낼 때 실행돼요(솔리디티(Solidity, 계약을 짜는 프로그래밍 언어)의 \texttt{emit} 문장이에요). 저장 공간에 값을 써넣는 일은 로그를 남기지 않아요. 누구나 검사할 수 있는 평판 시스템이 되려면, 블록체인 밖의 데이터베이스보다 로그를 골라 써야 해요. 로그는 블록체인 역사에서 영원히 남는 부분이고, 어떤 전체 노드(블록체인 기록을 모두 가진 컴퓨터)나 색인 제공자(블록체인 기록을 찾기 쉽게 정리해 주는 곳)도 따로 다시 계산해 볼 수 있기 때문이에요 \parencite{wood2025yellowpaper}.

송금은 늘 되풀이되지만, 허락은 나중 송금이 지나갈 수 있도록 한 번만 내줘요. 그래서 $G_{\text{transfer}}$가 더 빽빽한 그래프일 것이라고 봐요. 그리고 L2의 낮은 수수료가 그 차이를 더 벌릴 수도 있어요. 그렇게 생긴 그래프의 크기는 \ref{ch:results}장에 있어요.

\dissertationSubheading[sec:approval-security]{보안과 관련된 이벤트인 허락}

무제한 허락(\texttt{type(uint256).max})을 하면 주인은 다시 허락하지 않아도 돼요. 하지만 그만큼 안전을 내줘요. 허락받은 계정이 공격자에게 넘어가면, 주인이 허락을 다시 설정할 때까지 주인의 토큰을 하나도 남김없이 도둑맞을 수 있어요. 그래서 허락은 기록이 아니라 힘(무언가를 할 수 있는 능력)이에요. 허락을 가진 쪽은, 그 허락이 다 쓰이거나 거두어질 때까지 주인의 토큰을 옮길 수 있어요 \parencite{vogelsteller2015erc20}. 거래를 이용한 피싱은 바로 이 점을 노려요. 가짜 웹사이트가 사용자를 꾀어, 사기꾼이 사용자의 자산을 가져갈 수 있게 하는 거래에 서명하게 해요 \parencite{he2023txphishscope}. 엔도스랭크는 토큰 허락을, 한 스펜더가 그 토큰을 옮길 수 있는 힘을 보증하는 일로 풀이해요. 어떤 계정은 송금을 많이 받으면서도, 다른 누구를 대신해 돈을 옮겨도 된다는 허락은 한 번도 받지 못할 수 있어요.

\textcite{wang2022unlimited}에서는 이더리움의 ERC-20 허락 가운데 대부분이 무제한 허락이라는 것을 찾아냈어요. 그리고 무제한 허락이 지닌 위험을 수로 재었어요. 이 연구의 자료에서는 무제한 허락이 0보다 큰 마지막 허락의 절반 가까이를 차지해요(그 수는 부록~\ref{sec:data-tokens}에 있어요). 그리고 무제한 허락은 $2^{256}-1\approx1.16\times10^{77}$ 원래 단위(소수점 조정을 하지 않은 가장 작은 단위)의 금액으로 나타나요. 엔도스랭크의 금액 무게는 무제한 허락에 꽉 찬 무게 1을 줘요. 끝이 있는 허락(유한한 허락)은 중앙값(한가운데 값) 금액의 몇 배쯤 되어야 겨우 이 무게에 가까워져요. 한편 C-PR의 허락 층은 어떤 허락이든 달러로 정한 상한(이보다 큰 금액은 이 값으로 자름)까지만 세요(\ref{sub:coupled-graph}절). 그래서 이 층에서는 무제한 허락도 큰 유한한 허락보다 더 무겁지 않아요. 주인의 절반이 조금 넘는 수가 무제한 허락을 적어도 하나 가지고 있어요.

\dissertationSubheading[sec:awp-in-depth]{AWP 자세히 보기: 송금 그래프 기준선이 재는 것}

적응형 가중 페이지랭크(AWP)는 \textcite{do2023awp}에서 내놓은 방법이에요. 이 방법은 거래 그래프의 주소들에 순위를 매겨요. 거래 하나하나는 자기 화살표에 값을 더해요. 그 값은 거래 시간에 대한 로지스틱 함수에, 거래 금액에 대한 상한이 있는 함수를 곱한 것이에요. 그래서 어떤 거래도 화살표의 무게에 1보다 많이 더하지 않아요. 걷기는 각 주소에서, 그 주소의 활발함(활동 정도)에 비례해 다시 출발해요. 활발함은 그 주소가 보낸 상대 주소마다, 그 상대에게 보낸 가장 최근 거래의 시간 무게를 구해 모두 더한 값이에요. 그다음 순위가 얼마나 정확한지를 입차수와 받은 양의 합(들어온 금액을 모두 더한 값)에 견주어 검증해요. 이 둘은 한 지갑이 끄는 경제적 관심을 대신 재는 대리 지표로 쓰여요. 로지스틱 감쇠의 모양(식~\eqref{eq:logistic-decay})과, \ref{ch:methodology}장과 \ref{ch:results}장의 입차수/받은 양의 합 대리 지표는 AWP 논문에서 가져왔어요. 그 논문은 매개변수(계산에 넣는 정해진 값)의 값을 알려 주지 않아요. 여기서 쓴 값은 \ref{sub:awp-graph}절에 있어요. 이 방법은 모두 온체인에서 이루어지고 순위만을 바탕으로 해요. 그래서 블록체인 밖의 신분은 필요 없어요.

엔도스랭크는 페이지랭크 아래에서 쓰는 AWP의 무게 함수와, AWP의 순위 상관 점검을 그대로 지켜요. 바꾸는 것은 화살표의 뜻과, 걷기가 다시 출발하는 곳이에요. 엔도스랭크는 고르게 다시 출발하고, AWP는 활발하게 보내는 주소에서 다시 출발해요(\ref{sub:endorserank-graph}절). \ref{ch:results}장에서 우리는 다음을 시험해요. 허락 화살표가 다른 구성개념을 재고, 허락 점검과 잘 맞고, 계산이 더 싼 점수를 내는지 봐요. $t_1$에서 매긴 점수가 $t_2$까지 생긴 새 허락을 예측할 수 있는지도 봐요. 그리고 허락 화살표와 송금 화살표를 함께 걷는 걷기 하나가, 둘 가운데 하나만 걷는 걷기보다 계약에 순위를 더 잘 매기는지 봐요.

\dissertationSubheading[sec:method-comparison-table]{그래프 평판 방법 견주기}

표~\ref{tab:method-comparison-lit}에서는 엔도스랭크를 이 장에서 다룬 다른 방법들과 견주어요. 견주는 기준은 화살표의 뜻, 시간을 다루는 방식, 검증을 짜는 방식이에요.

\begin{table}[htbp]
\centering
\caption{고른 그래프 평판 방법들의 비교(개념 수준).}
\label{tab:method-comparison-lit}
\fitwidth{%
\begin{tabular}{p{0.18\linewidth}p{0.24\linewidth}p{0.24\linewidth}p{0.28\linewidth}}
\toprule
방법 & 화살표의 뜻 & 시간 모형 & 흔한 검증 방식 \\
\midrule
PageRank (웹) \parencite{page1999pagerank} & 하이퍼링크 & 없음 / 고정된 한 순간의 모습 & 검색 품질 \\
EigenTrust \parencite{kamvar2003eigentrust} & 이웃끼리 매긴 믿음 점수 & 되풀이해 모으기 & P2P 파일이 진짜인지 \\
TrustRank \parencite{gyongyi2004trustrank} & 좋은 씨앗에서 나가는 링크 & 고정 / 씨앗 사용 & 스팸 순위 낮추기 \\
이더리움의 PageRank\slash HodgeRank \parencite{do2023hodgerank} & 토큰 송금 & 고정된 기간 & 이 논문에서 재현하지 않음 \\
AWP \parencite{do2023awp} & 거래, 상한이 있는 금액 무게 & 화살표와 다시 출발의 로지스틱 감쇠 & 입차수 / 받은 양의 합 \\
RiskProp \parencite{lin2025riskprop} & 거래, 익명 벗기기 점수와 함께 & 위험의 퍼뜨림 & 라벨이 붙은 불법 계정 \\
다중 / 다층 PageRank \parencite{halu2013multiplex,dedomenico2015ranking} & 한 점 집합 위의 여러 관계 & 고정된 층 & 위쪽 순위 점의 다재다능함(여러 층에서 두루 중요한 정도) \\
\textcite{humanpassportnddocs} & 화살표 없음; 주소 주인에 대한 무게를 단 자격 증명(온체인과 오프체인) & 지금 가진 자격 증명 묶음 & 통과 기준; 알려진 사람 주소와 시빌 주소 \\
EndorseRank (이 연구) & 다섯 토큰의 마지막 허락(달러 금액), AWP의 무게 함수 & 마지막 허락마다의 로지스틱 감쇠; 고르게 다시 출발하기 & 허락 점검; $t_2$의 새 허락 \\
C-PR / S-PR (이 연구) & 송금과 섞거나 송금의 씨앗이 되는 허락; 달러 금액, 허락은 상한에서 자름 & 허락 스냅숏과 송금의 로지스틱 감쇠 & 홀드아웃 라벨; 미리 정한 규칙을 고친 점수에 다시 적용함 \\
\bottomrule
\end{tabular}
}
\end{table}

표~\ref{tab:method-comparison-lit}에서 보듯이, 이 논문의 주된 기여는 연산자(걷기를 정하는 계산 규칙) 쪽에 있어요. 이 논문은 두 관계를 걷기 하나, 곧 C-PR로 결합해요. 그리고 이 결합이 얼마나 보탬이 되는지 재려고, 씨앗을 쓰는 빼 보기 실험(한 부분을 바꿔 그 몫을 재는 실험)을 함께 해요. 또 이 결합을 기간 밖에서, 두 구성개념 모두에 견주어 시험해요. 이 시험은 이 연구가 자료를 뽑기 전에 정해 둔 규칙을 따라요. 주 분석의 고친 점수에는 이 규칙을 재분석(다시 하는 분석)으로 적용해요(\ref{sec:amount-revision}절과 부록~\ref{sec:map-registered}). 이 결합은 살펴본 방법들이 쓰지 않는 화살표 위에 서 있어요. 엔도스랭크의 기여는 마지막 허락을 평판 순위의 화살표로 쓰고, 그것을 기간 밖에서 평가한 데 있어요. 순위를 매기는 알고리즘이나, 허락을 주소 사이의 연결로 읽는 일은 엔도스랭크의 기여가 아니에요(\ref{sec:allowance-endorsement}절).

\dissertationSubheading[sec:literature-summary-table]{살펴본 연구 정리}

표~\ref{tab:literature-summary}에는 이 연구의 바탕이 되는 자료들을 모았어요. 이 논문이 가져다 쓰는 강점과, 이 논문이 풀려고 하는 한계도 함께 적었어요.

{\footnotesize
\setlength{\tabcolsep}{4pt}
\begin{longtable}{p{0.15\linewidth}p{0.20\linewidth}p{0.19\linewidth}p{0.19\linewidth}p{0.18\linewidth}}
\caption{살펴본 연구 정리: 방법, 강점, 한계, 이 연구와의 관계.}
\label{tab:literature-summary}\\
\toprule
자료(연도) & 방법 & 강점 & 한계 & 이 연구와의 관계 \\
\midrule
\endfirsthead
\multicolumn{5}{l}{\small 표~\thetable\ (이어짐)}\\
\toprule
자료(연도) & 방법 & 강점 & 한계 & 이 연구와의 관계 \\
\midrule
\endhead
\midrule
\multicolumn{5}{r}{\small 다음 쪽에 이어짐}\\
\endfoot
\bottomrule
\endlastfoot
\textcite{page1999pagerank}; \textcite{langville2006pagerank} & 하이퍼링크 그래프의 고유벡터 순위, 거듭제곱 반복으로 풀기 & 해석하기 쉬운 퍼뜨림; 반복 한 번에 $O(|E|)$; 감쇠 계수 덕분에 페론 벡터가 하나뿐임 & 화살표가 하이퍼링크만 뜻할 수 있음; 시간 개념이 없음 & 같은 풀이 프로그램과 복잡도 논증을, 견주는 두 방법에 바꾸지 않고 씀 \\
\textcite{kamvar2003eigentrust}; \textcite{gyongyi2004trustrank} & P2P 네트워크에서 믿음 모으기; 스팸 순위를 낮추려고 씨앗에서 퍼뜨리기 & 정규화와 씨앗 집합으로 부풀리기를 누름; 공격자를 분명하게 모형에 넣음 & 직접 매긴 평점이나 골라 둔 씨앗이 필요함; 온체인에 쓰인 적 없음 & 시빌 비용 논의와 고른 순간 이동 기준선의 바탕이 됨 \\
\textcite{do2023hodgerank} & 이더리움 거래 그래프에 PageRank와 HodgeRank를 써서 사회적 신용을 잼 & 송금 연결을 신용과 관련된 것으로 봄; 자료가 모두 공개됨 & 송금은 지불과 믿음을 뒤섞음; 허락의 뜻이 빠져 있음 & 앞선 송금 그래프 연구로, 이 논문에서는 AWP와 따로 다룸 \\
\textcite{do2023awp}, AWP & 상한이 있는 금액 무게, 로지스틱 최근성, 보낸 활동으로 무게를 단 활동 재시작을 쓴 거래 PageRank; 입차수/받은 양의 합과의 순위 상관을 검증으로 씀 & 순위로 검증해서 빚을 못 갚았다는 라벨이 필요 없음; 최근성을 분명하게 모형에 넣음 & 대리 지표를 그래프를 만든 바로 그 송금 자료로 계산함(기계적으로 생기는 일치); 빽빽한 그래프; 매개변수 값을 알리지 않음; 송금 화살표가 없는 계약은 모두 0에서 동점이 됨 & 기준선으로, 발표된 식을 바탕으로 다시 구현하고 $k$, $t_0$, $b$는 여기서 정함(\ref{sub:awp-graph}절); EndorseRank는 같은 무게 함수를 자기 기울기와 함께 씀 \\
\textcite{lin2025riskprop}, RiskProp & 계정 위험을 연결망으로 퍼뜨리고, 익명 벗기기 점수와 짝지음 & 주소 묶기와 분류로는 평가하지 못한 계정까지 모든 계정을 평가함; 라벨이 붙은 계정으로 점검함 & 불법 계정 라벨로 점검함; 믿음이나 허락이 아니라 사기를 겨냥함 & 거래를 따라 퍼뜨리면 계정 하나씩 보는 방법에 보탬이 됨을 보여 줌; 이 연구는 사기 라벨을 쓰지 않으므로 여기서 재현하지 않음 \\
\textcite{wu2022phishers} & 이더리움 거래 그래프의 연결망 임베딩으로 피싱 찾기 & 거래의 금액과 시간을 쓰는 임베딩; 앞선 임베딩 방법보다 효과가 좋음 & 라벨이 붙은 피싱 주소가 필요함; 라벨은 신용도가 아니라 피싱을 표시함 & 그래프 속 위치가 정보를 담는다는 전제를 뒷받침함; 경쟁하는 방법은 아님 \\
\textcite{hassija2020lending}; \textcite{kotb2023credit} & 전망 이론과 기계 학습을 쓴, 블록체인 자료 신용 점수 & 대출 결정을 바로 겨냥함 & 라벨이나 블록체인 밖의 특징이 필요함; 일부만 해석할 수 있음 & 대조되는 부류; EndorseRank는 검사할 수 있는 그래프 퍼뜨림 안에 머묾 \\
\textcite{bellini2020trust}; \textcite{almasoud2020reputation} & 블록체인 기반 믿음·평판 시스템 조사 & 설계 공간의 분류; 스마트 계약 위에 만든 평판 시스템 연구 서른 편 검토 & 평판을 서비스에 대한 사용자 평점에서 얻음; 권한 위임 화살표는 생각하지 않음; L2에서 잰 것이 없음 & 조사한 설계 가운데 마지막 허락 화살표를 쓰는 것이 없음 \\
\textcite{victor2020clustering}; \textcite{victor2021washtrading}; \textcite{cong2023washtrading} & 주소 묶기 어림 방법(그 가운데 하나는 ERC-20 허락으로 주소를 이음); 탈중앙 거래소(DEX)와 거래소에서 자전 거래 찾기 & 온체인 활동 가운데 가짜이거나 여러 주소에서 나온 몫을 잼 & 평판 점수를 내놓지 않음 & 허락을 쓴 가장 가까운 앞선 예로, 순위 화살표가 아니라 묶기 연결로 씀; 송금 기반 신호가 왜 부풀리기에 약한지 보여 줌 \\
\textcite{cronbach1955construct}; \textcite{messick1995validity} & 심리 측정의 구성 타당도(점수가 재려는 것을 정말 재는지) & 점수가 재도록 설계된 것과 바깥 기준을 구별함 & 블록체인에만 해당하는 것이 아님 & 같은 시기 점검을 구성개념 점검으로, 시간 홀드아웃을 기간 밖 증거로 다루는 바탕 \\
\textcite{humanpassportnddocs}; \textcite{siddarth2020watchmen}; \textcite{bordeianu2026sybil} & 주소 주인에 대한 무게를 단 자격 증명(온체인과 오프체인)을 통과 기준에 견줌; 사람 증명 방어 방법 검토 & 실제로 쓰이고 있음; 시빌 공격에 맞서 블록체인 밖의 증거를 더함 & 등급을 매긴 순서가 아니라 통과 여부만 가르는 문; 주소 사이 관계는 아무 역할이 없음; 시빌 전략이 계속 바뀜 & 실제로 믿음을 어떻게 세우는지, 그리고 EndorseRank가 주지 못하는 것(사람임을 보증하지 않음)을 보여 줌 \\
\textcite{efron1979bootstrap}; \textcite{diciccio1996bootstrap} & 부트스트랩 다시 뽑기와 부트스트랩 기반 신뢰구간 & 분포 가정 없이 순위 통계의 불확실성을 구함; 짝지은 대비 & 단위들이 서로 바꿔도 되는 것(교환 가능)이어야 함 & \ref{ch:results}장에서 보고하는 95\% 구간과 $\Delta\tau$ 대비의 출처 \\
\end{longtable}
}

\dissertationSubheading[sec:research-gap]{연구의 빈틈}

표~\ref{tab:literature-summary}와 \ref{sec:related-measurement}절에서 다섯 가지 빈틈이 드러나요. 첫째, 여기서 다룬 온체인 평판 방법들은 모두 송금이나 직접 매긴 평점으로 그래프를 만들어요(페이지랭크/HodgeRank \parencite{do2023hodgerank}, AWP \parencite{do2023awp}, RiskProp \parencite{lin2025riskprop}, 그리고 조사 연구들에 실린 설계들 \parencite{bellini2020trust}). 그리고 디파이에서 쓰이는 증명 도구는, 주소 사이 관계의 그래프 없이 주소를 가진 사람에게 점수를 매겨요 \parencite{humanpassportnddocs}. 어느 것도 ERC-20 허락을 쓰지 않아요. 허락은 쓸 권한을 맡긴 일의 온체인 기록이고, 모든 ERC-20 토큰에 들어 있어요. 허락은 주소 묶기와 보안 분석에 쓰여 왔어요(\ref{sec:allowance-endorsement}절과 \ref{sec:approval-security}절). 하지만 여기서 살펴본 연구에서는 허락이 평판 순위의 화살표로 쓰인 적이 없어요. 마지막 허락으로 만든 보증 그래프가 남다르고 뜻있는 순서를 내는지는 아직 알려지지 않았어요(O1과 RQ1). 둘째 빈틈은 검증에 있어요. 송금 그래프 기준선은 자기 점수를, 그래프를 만드는 데 쓴 바로 그 송금 자료로 계산한 입차수와 받은 양의 합에 견주어 점검해요. 그래서 어느 정도 맞는 것은 피할 수 없어요. 이 절차를 허락 구성개념과 그 짝인 송금 구성개념에 쓰고, 같은 주소 묶음에서 신뢰구간과 짝지은 차이 검사까지 함께 보고한 연구는 없어요(O1과 RQ1). 셋째, 여기서 살펴본 송금 그래프 연구들은 정해진 절차 아래에서 잰 실행 시간을 화살표 수, 메모리 사용량, 반복 횟수와 함께 내놓지 않아요. 그래서 효율이 좋다는 주장이 있어도, 그것이 구현 때문이 아니라 희소성 때문이라고 짚어 내기 어려워요(O2와 RQ2). 넷째 빈틈은 방법에 관한 것이고, 시간과 관련이 있어요. 같은 기간의 그래프 통계를 검증으로 내세우는 일이 많아요. 하지만 그런 점수로는 앞으로의 행동을 예측한다는 것을 보일 수 없어요. 그래서 여기서 쓰는 기간 밖 점검은 그 대신 시간 홀드아웃이에요. 구성개념마다 앞으로 일어날 일을 세는 라벨을 두고, 동결 때의 원시 차수를 기준선으로 삼아요(O3과 RQ3). 마지막으로, 다중 페이지랭크 \parencite{halu2013multiplex}와 다층 순위 매기기 \parencite{dedomenico2015ranking}는 사회 연결망과, 여러 종류의 관계가 있는 다른 연결망에 쓰여 왔어요. 하지만 여기서 살펴본 연구에서는 한 주소 묶음 위의 두 온체인 관계에 쓰인 적이 없어요. 결과를 보기 전에 정해 둔 규칙 아래에서, 권한 위임 층과 지불 층을 합치면 두 층을 모두 이기는지, 또는 지불 층 하나만 쓸 때보다 나아지는지 보고한 연구는 없어요(O4와 RQ4, O5와 RQ5). 이 논문의 방법인 C-PR은 이 다섯째 빈틈을 다뤄요. 앞의 네 빈틈은 C-PR에 견줄 기준 걷기들과 평가 방법을 마련해 줘요.

\dissertationSubheading[sec:open-problems]{이 연구를 하게 만든 풀리지 않은 문제}

앞선 연구들은 온체인 평판에서 몇 가지 문제를 풀지 못한 채 남겨 두었어요. 이 문제들 때문에 이 연구를 하게 되었어요.

\begin{enumerate}
\item 화살표의 뜻: 신분이 가명일 때 \parencite{lin2024whoiswho}, 어떤 온체인 관계가 믿음이나 권한 위임을 가장 잘 담아낼까요 \parencite{packin2024credit}?
\item 채무 불이행 기록 없이 검증하기: 대출을 갚지 못했다는 라벨이 드물 때, 무엇을 바탕으로 방법들을 견줄 수 있을까요 \parencite{do2023hodgerank}?
\item 계산 비용: 그래프가 커져도, 평판 점수를 사람이 기다릴 만한 짧은 시간 안에 계속 새로 고칠 수 있을까요 \parencite{langville2006pagerank}?
\item 시간으로 나눈 점검: 무엇이 있어야 평가가 같은 기간의 그래프 통계를 기간 밖 증거로 잘못 여기지 않을까요 \parencite{cronbach1955construct}?
\item 공격에 버티는 힘: 신분을 싸게 만들 수 있을 때, 평판을 꾸며 내는 데 비용이 얼마나 들까요 \parencite{douceur2002sybil}?
\end{enumerate}

이 논문은 (1)--(4)를 허락 화살표와 그 화살표를 송금과 결합하는 걷기, 같은 기간 점검, 실행 시간 벤치마크(속도 재기), 시간 홀드아웃으로 다뤄요. 문제 (5)는 \ref{ch:discussion}장에서 모형 수준으로 분석해요. 실제 자료로 공격을 온전히 평가하는 일은 앞으로의 연구로 남겨 둬요.

\ref{ch:methodology}장에서는 여기서 밝힌 빈틈을 실제로 해 볼 수 있는 설계로 바꿔요.

